% !PRINTED 640
\weberpage{640}
%<<CONT
and consequently $\frac{1}{2}(\alpha \pm \beta)$ cannot be an integer, because the
absolute value of the norm of this number is less than 1. Now
if $\omega_1, \omega_2, \dots, \omega_n$ is a basis of $\mathfrak{o}$, and
\[
\begin{aligned}
\alpha &= a_1\omega_1 + a_2\omega_2 + \dots + a_n\omega_n \\
\beta &= b_1\omega_1 + b_2\omega_2 + \dots + b_n\omega_n,
\end{aligned}
\]
then the integral rational numbers $a, b$ cannot satisfy the congruences
\[ a_1 \equiv b_1,\ a_2 \equiv b_2,\ \dots,\ a_n \equiv b_n \pmod{2} \]
since otherwise $\frac{1}{2}(\alpha - \beta)$ would be an integer.

If, then, we distribute all the numbers $\alpha$ into $2^n$ compartments,
by throwing into one compartment all the numbers in which
$a_1, a_2, \dots, a_n$ leave the same residues (0 or 1) with respect to the modulus 2,
then in each of these compartments at most 2 numbers,
namely $\alpha$ and $-\alpha$, can occur; and if we further exclude
$\alpha = 0$, then in one of these compartments, namely that in which
the $a_1, a_2, \dots, a_n$ are all even, there is in fact no number $\alpha$ at all. The number
of all possible numbers $\alpha$ is therefore finite, and at most
$= 2^{n+1} - 2$.

If now the absolute value of $\alpha = 1$, then the same holds
of all the powers of $\alpha$. Consequently, in the unbounded
sequence
\[ 1,\ \alpha,\ \alpha^2,\ \alpha^3,\ \dots \]
the same number must necessarily recur a second time, and hence
$\alpha^h = \alpha^k$ and $k > h$. But then
\[ \alpha^{k-h} = 1, \]
that is, $\alpha$ is a root of unity, as was to be proved.

We further add the remark that if a number $\alpha$
of a field $\Omega$ is a root of unity of degree $m$, then all
the numbers conjugate to $\alpha$ are also roots of unity of the same degree.
For in the rational equation $\alpha^m - 1 = 0$, $\alpha$ can be
replaced by each of the conjugate values.

The number of roots of unity that are contained in a field $\Omega$
can only ever be a finite one, because every
number in $\Omega$ satisfies a rational equation whose degree is
at most equal to the degree of the field. The degree of the roots of unity
in $\Omega$ therefore cannot exceed a finite value.
